% Mathematics Bootcamp --- Day 1 Exercises, with Solutions
% Reviewed copy: logical scope, proof hierarchy, and blackboard readability.
% Exercise statements, numbering, notation, and proof strategies are retained.
% See review_notes.md for the substantive corrections and scope of the review.
% Mas-Colell, Whinston and Green, Chapter 1, Sections 1.B and 1.C.
%
% Self-contained: this file needs no other file in the repository.
% Compile with pdflatex, twice, or with latexmk -pdf.

\documentclass[12pt]{article}

\usepackage[margin=1in]{geometry}

\usepackage{
  amsmath,
  amssymb,
  amsfonts,
  amsthm,
}

\usepackage[T1]{fontenc}
\usepackage{lmodern}

\usepackage{setspace}
\usepackage{float}

% Added to the template: the roman/alpha labels in the exercise statements
% need it. source/proofs/main.tex already loads it too.
\usepackage{enumitem}
\usepackage{needspace}

% Keep each current proof together. The longproof environment allows page breaks.
\raggedbottom
\clubpenalty=10000
\widowpenalty=10000
\displaywidowpenalty=10000

\onehalfspacing

% ---------------------------------------------------------------------------
% Structural elements: exercise headings, explicit goals, and nested proof
% blocks. All logical indentation is controlled by \plevel below.
% ---------------------------------------------------------------------------

% \exercise{number}{MWG grade}{one-line description}
\newcommand{\exercise}[3]{%
  \par\penalty-250\bigskip\Needspace{8\baselineskip}
  \noindent\textbf{\large Exercise #1}\quad\textnormal{(grade #2)}%
  \par\nobreak\smallskip
  \noindent\textit{#3}%
  \par\nobreak\smallskip
  \noindent\rule{\linewidth}{0.4pt}%
  \par\nobreak\smallskip}

% The statement, verbatim from the book.
\newenvironment{statement}{\par\medskip}{\par\medskip}

% What the proof is obliged to produce, written before any of it is written.
\newenvironment{nts}
  {\par\medskip\Needspace{3\baselineskip}%
   \noindent\textbf{Need to Show:}\par\nobreak\smallskip}
  {\par\medskip}

% One level of indentation. Nest to indent further: the structure of the
% argument is carried by the indentation, as in the Day 1 slides.
\newlength{\plevel}\setlength{\plevel}{2.25em}
\makeatletter
\newenvironment{pstep}
  {\par\nobreak\begin{list}{}{%
     \setlength{\leftmargin}{\plevel}%
     \setlength{\rightmargin}{0pt}%
     \@beginparpenalty=10000\relax
     \@endparpenalty=10000\relax
     \setlength{\labelwidth}{0pt}\setlength{\labelsep}{0pt}%
     \setlength{\itemindent}{0pt}\setlength{\listparindent}{0pt}%
     \setlength{\topsep}{0pt}\setlength{\partopsep}{0pt}%
     \setlength{\parsep}{0pt}\setlength{\itemsep}{0pt}}%
   \item\relax}
  {\end{list}}
\makeatother

% A proof heading is not a logical premise: put it on its own line.
% Suppress ordinary paragraph indents inside proofs, so that only pstep
% creates a new logical level. Preserve amsthm's end-of-proof mechanism.
\makeatletter
\newif\ifboardproofboxed
\boardproofboxedtrue
\renewenvironment{proof}[1][\proofname]
  {\par\addvspace{.7\baselineskip}\Needspace{5\baselineskip}%
   \ifboardproofboxed\noindent\begin{minipage}{\linewidth}\fi
   \pushQED{\qed}\normalfont
   \topsep0pt\relax
   \trivlist\item\relax
   \setlength{\parindent}{0pt}\setlength{\listparindent}{0pt}%
   \setlength{\parskip}{0pt}%
   \textit{#1.}\par\nobreak}
  {\popQED\endtrivlist\@endpefalse
   \ifboardproofboxed\end{minipage}\par\fi}
% Use longproof for a future argument that is too long for one page.
\newenvironment{longproof}[1][\proofname]
  {\boardproofboxedfalse\begin{proof}[#1]}
  {\end{proof}}
\makeatother

% Right-flushed justification for a line of a proof. Ends the line.
\newcommand{\by}[1]{%
  \unskip\nobreak\hfil\penalty50\hskip1em\null\nobreak\hfil
  \textnormal{[#1]}%
  {\parfillskip=0pt \finalhyphendemerits=0 \par}}

\newenvironment{remark}
  {\par\medskip\noindent\textit{Remark.}\enspace}
  {\par\medskip}

% Contradiction. The single arrow is the one the Day 1 slides use for
% implication, so this reads as "then, and then back the other way".
\newcommand{\contra}{\ensuremath{\rightarrow\leftarrow}}

\title{\textbf{Mathematics Bootcamp --- Day 1 Exercises, with Solutions}\\[.4em]
       \large Mas-Colell, Whinston and Green, Chapter 1}
\author{Miguel V\'azquez-Carrero}
\date{Preference and Choice: Sections 1.B and 1.C\\[.3em]
      \normalsize Compilation Date: \today}

\begin{document}

\maketitle

\section*{Notation}

Throughout, $X$ is a set of alternatives and $\succeq$ is a binary relation on
$X$, read ``at least as good as''. From it we derive
\[
x\succ y \iff x\succeq y \ \text{ and not } y\succeq x,
\qquad
x\sim y \iff x\succeq y \ \text{ and } y\succeq x .
\]
The relation $\succeq$ is \emph{rational} if it is
\emph{complete} (for all $x,y\in X$, $x\succeq y$ or $y\succeq x$)
and \emph{transitive} (for all $x,y,z\in X$, if $x\succeq y$ and $y\succeq z$
then $x\succeq z$). A function $u:X\to\mathbb R$ \emph{represents} $\succeq$ if
\[
x\succeq y \iff u(x)\geq u(y)\qquad\text{for all }x,y\in X .
\]

A \emph{choice structure} is a pair $(\mathbb B,C(\cdot))$: $\mathbb B$ is a
family of nonempty subsets of $X$, the observed menus, and $C$ assigns to each
$B\in\mathbb B$ a nonempty set $C(B)\subseteq B$. MWG writes $\mathcal B$ for
the menu domain; we keep the $\mathbb B$ of the slides.

\begin{quote}
  \textbf{Definition 1.C.1 (Weak Axiom).}
  If for some $B\in\mathbb B$ with $x,y\in B$ we have $x\in C(B)$, then for any
  $B^{\prime}\in\mathbb B$ with $x,y\in B^{\prime}$ and $y\in C(B^{\prime})$, we
  must also have $x\in C(B^{\prime})$.

  \medskip
  \textbf{Definition 1.C.2 (Revealed preference).}
  $x\succeq^{*}y \iff$ there is some $B\in\mathbb B$ such that $x,y\in B$ and
  $x\in C(B)$.
\end{quote}

We abbreviate the weak axiom as WARP (Weak Axiom of Revealed Preference).

Exercises 1.B.1 and 1.B.2 ask for the proof of the following, which MWG states
and does not prove.

\begin{quote}
  \textbf{Proposition 1.B.1.} If $\succeq$ is rational then:
  \begin{enumerate}[label=(\roman*),leftmargin=2.4em,topsep=.3em,itemsep=.2em]
    \item $\succ$ is both \emph{irreflexive} ($x\succ x$ never holds) and
    \emph{transitive} (if $x\succ y$ and $y\succ z$, then $x\succ z$).
    \item $\sim$ is \emph{reflexive} ($x\sim x$ for all $x$), \emph{transitive}
    (if $x\sim y$ and $y\sim z$, then $x\sim z$), and \emph{symmetric}
    (if $x\sim y$, then $y\sim x$).
    \item if $x\succ y\succeq z$, then $x\succ z$.
  \end{enumerate}
\end{quote}

\section*{How to write these up}

Write the skeleton first, then the body. For a universally quantified
implication, the structure is:
\begin{quote}
  Let $x$ be arbitrary.
  \begin{pstep}
    Suppose $P(x)$.
    \begin{pstep}
      [Prove $Q(x)$ using $P(x)$ and the standing hypotheses.]
    \end{pstep}
    Thus $P(x)\Rightarrow Q(x)$.
  \end{pstep}
  Thus, for all $x$, $P(x)\Rightarrow Q(x)$.
\end{quote}

Indentation records logical dependence. A temporary assumption is available
inside its block and every sub-block. Previously established facts remain
available on later lines at the same level and in deeper blocks. Returning to
the level of a temporary assumption closes its block: carry out the justified
conclusion, not the assumption itself. For example, after deriving a
contradiction from $R$, return one level and write ``Thus not $R$.''

A \emph{claim} is a goal to prove, not an extra hypothesis. Indent its proof;
once proved, the claim is available in the surrounding argument. Put parallel
claims, cases, and directions at the same level. Give each case its own block,
and return to their common level before combining the cases.

Before the proof, state the \textbf{Need to Show}. Unfold the definitions
relevant to the conclusion until the required implications, conjunctions,
witnesses, or counterexamples are explicit. Keep established notation and
lemmas when they make the target clearer. A proof is finished when every
stated obligation has been met.

\newpage
\section*{Preferences (Section 1.B)}

% ===========================================================================
\exercise{1.B.1}{B}{Property (iii) of Proposition 1.B.1}

\begin{statement}
  Prove property (iii) of Proposition 1.B.1.
\end{statement}

\begin{nts}
  For all $x,y,z\in X$: if $x\succ y$ and $y\succeq z$, then $x\succ z$.
  Unfolding $\succ$, the hypothesis is
  \[
  \underbrace{x\succeq y \ \text{ and not } y\succeq x}_{x\,\succ\,y},
  \qquad y\succeq z,
  \]
  and the conclusion $x\succ z$ is two separate claims:
  \begin{enumerate}[label=(\alph*),leftmargin=2.4em,topsep=.2em,itemsep=.1em]
    \item $x\succeq z$;
    \item not $z\succeq x$.
  \end{enumerate}
\end{nts}

\begin{proof}
  Let $x,y,z\in X$ be arbitrary.
  \begin{pstep}
    Suppose $x\succ y$ and $y\succeq z$.
    \begin{pstep}
      \emph{Claim (a): $x\succeq z$.}
      \begin{pstep}
        $x\succeq y$. \by{definition of $\succ$}

        $x\succeq y$ and $y\succeq z$, so $x\succeq z$. \by{transitivity}
      \end{pstep}
      \emph{Claim (b): not $z\succeq x$.}
      \begin{pstep}
        Suppose, for contradiction, $z\succeq x$.
        \begin{pstep}
          $y\succeq z$ and $z\succeq x$, so $y\succeq x$. \by{transitivity}

          But $x\succ y$ says not $y\succeq x$. \contra
        \end{pstep}
        Thus not $z\succeq x$.
      \end{pstep}
      By (a) and (b), $x\succeq z$ and not $z\succeq x$, that is, $x\succ z$.
    \end{pstep}
    Thus, if $x\succ y$ and $y\succeq z$, then $x\succ z$.
  \end{pstep}
  Thus, for all $x,y,z\in X$, if $x\succ y\succeq z$ then $x\succ z$.
\end{proof}

\begin{remark}
  Only transitivity is used; completeness is never invoked. The two claims (a)
  and (b) are proved by different methods --- (a) directly, (b) by contradiction
  --- which is the usual shape whenever the conclusion is a strict preference:
  one half builds a chain, the other half rules out the reverse chain.
\end{remark}

% ===========================================================================
\exercise{1.B.2}{A}{Properties (i) and (ii) of Proposition 1.B.1}

\begin{statement}
  Prove properties (i) and (ii) of Proposition 1.B.1.
\end{statement}

\begin{nts}
  Five separate claims. For all $x,y,z\in X$:
  \begin{enumerate}[label=(\alph*),leftmargin=2.4em,topsep=.2em,itemsep=.1em]
    \item $\succ$ is irreflexive: not $x\succ x$;
    \item $\succ$ is transitive: if $x\succ y$ and $y\succ z$, then $x\succ z$;
    \item $\sim$ is reflexive: $x\sim x$;
    \item $\sim$ is transitive: if $x\sim y$ and $y\sim z$, then $x\sim z$;
    \item $\sim$ is symmetric: if $x\sim y$, then $y\sim x$.
  \end{enumerate}
  Claim (c) uses $x\succeq x$, which follows from completeness and is proved
  first as Step 0. Claim (a) follows directly from the definition of $\succ$;
  it needs neither reflexivity nor rationality.
\end{nts}

\begin{proof}[Step 0: $\succeq$ is reflexive]
  Let $x\in X$ be arbitrary.
  \begin{pstep}
    Completeness applied to the pair $(x,x)$ gives: $x\succeq x$ or $x\succeq x$.

    Either disjunct is $x\succeq x$.
  \end{pstep}
  Thus $x\succeq x$ for every $x\in X$.
\end{proof}

\begin{proof}[Proof of (a), irreflexivity of $\succ$]
  Let $x\in X$ be arbitrary.
  \begin{pstep}
    Suppose, for contradiction, $x\succ x$.
    \begin{pstep}
      Then $x\succeq x$ and not $x\succeq x$. \by{definition of $\succ$}

      These two conclusions contradict each other. \contra
    \end{pstep}
    Thus not $x\succ x$.
  \end{pstep}
  Thus $\succ$ is irreflexive.
\end{proof}

\Needspace{11\baselineskip}
\begin{proof}[Proof of (b), transitivity of $\succ$]
  Let $x,y,z\in X$ be arbitrary.
  \begin{pstep}
    Suppose $x\succ y$ and $y\succ z$.
    \begin{pstep}
      $y\succeq z$. \by{definition of $\succ$}

      So $x\succ y$ and $y\succeq z$, hence $x\succ z$.
      \by{Proposition 1.B.1(iii), Exercise 1.B.1}
    \end{pstep}
    Thus, if $x\succ y$ and $y\succ z$, then $x\succ z$.
  \end{pstep}
  Thus $\succ$ is transitive.
\end{proof}

\begin{proof}[Proof of (c), reflexivity of $\sim$]
  Let $x\in X$ be arbitrary.
  \begin{pstep}
    $x\succeq x$. \by{Step 0}

    So $x\succeq x$ and $x\succeq x$, that is, $x\sim x$. \by{definition of $\sim$}
  \end{pstep}
  Thus $\sim$ is reflexive.
\end{proof}

\begin{proof}[Proof of (d), transitivity of $\sim$]
  Let $x,y,z\in X$ be arbitrary.
  \begin{pstep}
    Suppose $x\sim y$ and $y\sim z$.
    \begin{pstep}
      Unfolding: $x\succeq y$, $y\succeq x$, $y\succeq z$, $z\succeq y$.

      $x\succeq y$ and $y\succeq z$, so $x\succeq z$. \by{transitivity}

      $z\succeq y$ and $y\succeq x$, so $z\succeq x$. \by{transitivity}

      So $x\succeq z$ and $z\succeq x$, that is, $x\sim z$.
    \end{pstep}
    Thus, if $x\sim y$ and $y\sim z$, then $x\sim z$.
  \end{pstep}
  Thus $\sim$ is transitive.
\end{proof}

\begin{proof}[Proof of (e), symmetry of $\sim$]
  Let $x,y\in X$ be arbitrary.
  \begin{pstep}
    Suppose $x\sim y$.
    \begin{pstep}
      Unfolding: $x\succeq y$ and $y\succeq x$.

      The same two facts, in the other order, are $y\succeq x$ and $x\succeq y$,
      that is, $y\sim x$. \by{definition of $\sim$}
    \end{pstep}
    Thus, if $x\sim y$ then $y\sim x$.
  \end{pstep}
  Thus $\sim$ is symmetric.
\end{proof}

\begin{remark}
  Note the order of dependence in this presentation: the proof of (b) uses
  property (iii), established in Exercise 1.B.1. This is a convenient order of
  proof, not a logically necessary order: (b) could also be proved directly.

  Note also that reflexivity of $\succeq$ is a consequence of completeness, not a
  separate assumption --- completeness quantifies over all pairs $x,y\in X$,
  including the pair $(x,x)$. This is why the definition of rationality has two
  clauses and not three. In contrast, irreflexivity of $\succ$ and symmetry
  of $\sim$ follow from their definitions for any underlying relation.
\end{remark}

% ===========================================================================
\exercise{1.B.3}{B}{A strictly increasing transformation of a utility function}

\begin{statement}
  Show that if $f:\mathbb R\to\mathbb R$ is a strictly increasing function and
  $u:X\to\mathbb R$ is a utility function representing preference relation
  $\succeq$, then the function $v:X\to\mathbb R$ defined by $v(x)=f(u(x))$ is
  also a utility function representing preference relation $\succeq$.
\end{statement}

\begin{nts}
  For all $x,y\in X$: $x\succeq y \iff v(x)\geq v(y)$.
  \medskip

  \noindent
  We are given $x\succeq y\iff u(x)\geq u(y)$, so it is enough to show
  $u(x)\geq u(y)\iff f(u(x))\geq f(u(y))$. That is a statement about $f$ alone:
  for all $a,b\in\mathbb R$, $a\geq b\iff f(a)\geq f(b)$. It is proved first, as
  a lemma. Note what is available and what is wanted: strictly increasing is an
  implication about $>$, and the claim is an equivalence about $\geq$.
\end{nts}

\begin{proof}[Lemma]
  Let $a,b\in\mathbb R$ be arbitrary.
  \begin{pstep}
    $(\Rightarrow)$ Suppose $a\geq b$.
    \begin{pstep}
      \emph{Case 1: $a=b$.}
      \begin{pstep}
        Then $f(a)=f(b)$, so $f(a)\geq f(b)$.
      \end{pstep}
      \emph{Case 2: $a>b$.}
      \begin{pstep}
        Then $f(a)>f(b)$, so $f(a)\geq f(b)$.
        \by{$f$ strictly increasing}
      \end{pstep}
      Since $a\geq b$, these cases are exhaustive. Both give $f(a)\geq f(b)$.
    \end{pstep}
    Thus $a\geq b\Rightarrow f(a)\geq f(b)$.

    $(\Leftarrow)$ We prove the contrapositive. Suppose not $a\geq b$.
    \begin{pstep}
      Then $b>a$. \by{trichotomy in $\mathbb R$}

      So $f(b)>f(a)$, that is, not $f(a)\geq f(b)$.
      \by{$f$ strictly increasing}
    \end{pstep}
    Thus not $a\geq b$ implies not $f(a)\geq f(b)$.

    By contraposition, $f(a)\geq f(b)\Rightarrow a\geq b$.
  \end{pstep}
  Thus, for all $a,b\in\mathbb R$, $a\geq b\iff f(a)\geq f(b)$.
\end{proof}

\Needspace{11\baselineskip}
\begin{proof}
  Let $x,y\in X$ be arbitrary.
  \begin{pstep}
    $x\succeq y \iff u(x)\geq u(y)$. \by{$u$ represents $\succeq$}

    $u(x)\geq u(y) \iff f(u(x))\geq f(u(y))$. \by{Lemma, with $a=u(x)$, $b=u(y)$}

    $f(u(x))\geq f(u(y)) \iff v(x)\geq v(y)$. \by{definition of $v$}

    Chaining the three equivalences, $x\succeq y\iff v(x)\geq v(y)$.
  \end{pstep}
  Thus $v$ represents $\succeq$.
\end{proof}

\begin{remark}
  The whole content is in the Lemma, and the Lemma is where strictness is used.
  If $f$ were merely nondecreasing, the direction $(\Leftarrow)$ can fail: take $f$
  constant, and $v$ represents only the trivial preference in which everything is
  indifferent to everything else. This is the formal content of the claim that
  utility is \emph{ordinal}. The numbers carry no information beyond the order
  they induce, so any relabelling that preserves that order represents the same
  preference.
\end{remark}

% ===========================================================================
\exercise{1.B.4}{A}{A sufficient condition for representation}

\begin{statement}
  Consider a rational preference relation $\succeq$. Show that if $u(x)=u(y)$
  implies $x\sim y$ and if $u(x)>u(y)$ implies $x\succ y$, then $u(\cdot)$ is a
  utility function representing $\succeq$.
\end{statement}

\begin{nts}
  For all $x,y\in X$: $x\succeq y\iff u(x)\geq u(y)$.
  \medskip

  \noindent
  The two hypotheses run from numbers to preference. Representation requires both
  directions, so the work is to reverse them. The device is trichotomy: fix
  $x,y$, and observe that each side of the equivalence splits into exactly three
  mutually exclusive and exhaustive cases,
  \[
  \begin{array}{lccc}
    \text{on the right:} & u(x)>u(y), & u(x)=u(y), & u(y)>u(x);\\[.35em]
    \text{on the left:}  & x\succ y,  & x\sim y,   & y\succ x.
  \end{array}
  \]
  Each numerical case implies the corresponding preference case. If one
  preference case holds, exclusivity rules out the other two numerical cases;
  exhaustiveness then forces the corresponding numerical case. Step 3 makes
  these reverse implications explicit.
\end{nts}

\begin{proof}[Step 1: the three numerical cases]
  Let $x,y\in X$ be arbitrary.
  \begin{pstep}
    Exactly one of $u(x)>u(y)$, $u(x)=u(y)$, and $u(y)>u(x)$ holds.
    \by{trichotomy in $\mathbb R$}
  \end{pstep}
  Thus the numerical cases are mutually exclusive and exhaustive for every pair.
\end{proof}

\begin{proof}[Step 2: the three preference cases]
  Let $x,y\in X$ be arbitrary.
  \begin{pstep}
    \emph{At least one preference case holds.}
    \begin{pstep}
      $x\succeq y$ or $y\succeq x$. \by{completeness}

      \emph{Case 1: both $x\succeq y$ and $y\succeq x$.}
      \begin{pstep}
        Then $x\sim y$. \by{definition of $\sim$}
      \end{pstep}
      \emph{Case 2: $x\succeq y$ but not $y\succeq x$.}
      \begin{pstep}
        Then $x\succ y$. \by{definition of $\succ$}
      \end{pstep}
      \emph{Case 3: $y\succeq x$ but not $x\succeq y$.}
      \begin{pstep}
        Then $y\succ x$. \by{definition of $\succ$}
      \end{pstep}
      Completeness makes these cases exhaustive, so at least one preference case holds.
    \end{pstep}
    \emph{At most one preference case holds.}
    \begin{pstep}
      Suppose $x\succ y$ and $x\sim y$ both hold.
      \begin{pstep}
        The first gives not $y\succeq x$; the second gives $y\succeq x$. \contra
      \end{pstep}
      Thus $x\succ y$ and $x\sim y$ cannot both hold.

      Suppose $x\succ y$ and $y\succ x$ both hold.
      \begin{pstep}
        The first gives $x\succeq y$; the second gives not $x\succeq y$. \contra
      \end{pstep}
      Thus $x\succ y$ and $y\succ x$ cannot both hold.

      Suppose $x\sim y$ and $y\succ x$ both hold.
      \begin{pstep}
        The first gives $x\succeq y$; the second gives not $x\succeq y$. \contra
      \end{pstep}
      Thus $x\sim y$ and $y\succ x$ cannot both hold.

      These are all pairs of distinct preference cases, so at most one holds.
    \end{pstep}
    Combining the two claims, exactly one preference case holds.
  \end{pstep}
  Thus, for all $x,y\in X$, exactly one of $x\succ y$, $x\sim y$, $y\succ x$ holds.
\end{proof}

\begin{proof}[Step 3: the hypotheses reverse]
  Let $x,y\in X$ be arbitrary.
  \begin{pstep}
    The hypotheses give
    \[
    \begin{aligned}
      u(x)>u(y)&\Rightarrow x\succ y,\\
      u(x)=u(y)&\Rightarrow x\sim y,\\
      u(y)>u(x)&\Rightarrow y\succ x.
    \end{aligned}
    \]
    The third is the first with $x$ and $y$ exchanged.

    Suppose $x\succ y$.
    \begin{pstep}
      Suppose, for contradiction, $u(x)=u(y)$.
      \begin{pstep}
        Then $x\sim y$ by hypothesis, contrary to $x\succ y$. \contra
        \by{exclusivity in Step 2}
      \end{pstep}
      Thus $u(x)\neq u(y)$.

      Suppose, for contradiction, $u(y)>u(x)$.
      \begin{pstep}
        Then $y\succ x$ by hypothesis, contrary to $x\succ y$. \contra
        \by{exclusivity in Step 2}
      \end{pstep}
      Thus not $u(y)>u(x)$.

      The only remaining numerical case is $u(x)>u(y)$. \by{Step 1}
    \end{pstep}
    Thus $x\succ y\Rightarrow u(x)>u(y)$.

    Suppose $x\sim y$.
    \begin{pstep}
      Suppose, for contradiction, $u(x)>u(y)$.
      \begin{pstep}
        Then $x\succ y$ by hypothesis, contrary to $x\sim y$. \contra
        \by{exclusivity in Step 2}
      \end{pstep}
      Thus not $u(x)>u(y)$.

      Suppose, for contradiction, $u(y)>u(x)$.
      \begin{pstep}
        Then $y\succ x$ by hypothesis, contrary to $x\sim y$. \contra
        \by{exclusivity in Step 2}
      \end{pstep}
      Thus not $u(y)>u(x)$.

      The only remaining numerical case is $u(x)=u(y)$. \by{Step 1}
    \end{pstep}
    Thus $x\sim y\Rightarrow u(x)=u(y)$.

    The same strict-preference argument, with $x$ and $y$ exchanged, gives
    $y\succ x\Rightarrow u(y)>u(x)$.
  \end{pstep}
  Thus all three implications are equivalences for every $x,y\in X$.
\end{proof}

\begin{proof}[Step 4: conclusion]
  Let $x,y\in X$ be arbitrary.
  \begin{pstep}
    $x\succeq y\iff (x\succ y\text{ or }x\sim y)$.
    \by{definitions of $\succ$ and $\sim$}

    $(x\succ y\text{ or }x\sim y)\iff (u(x)>u(y)\text{ or }u(x)=u(y))$.
    \by{Step 3}

    $(u(x)>u(y)\text{ or }u(x)=u(y))\iff u(x)\geq u(y)$.
    \by{definition of $\geq$}

    Chaining these equivalences gives $x\succeq y\iff u(x)\geq u(y)$.
  \end{pstep}
  Thus the equivalence holds for all $x,y\in X$, so $u$ represents $\succeq$.
\end{proof}

\begin{remark}
  This proof never uses transitivity and invokes completeness only in Step 2.
  In fact, completeness also follows from the hypotheses: trichotomy of
  $u(x),u(y)$ forces $x\succ y$, $x\sim y$, or $y\succ x$, and hence at least
  one weak comparison. The representation then gives transitivity through the
  order on $\mathbb R$. The rationality assumption is therefore stronger than
  needed for this particular exercise.

  The first line of Step 4 also deserves a pause. It says $x\succeq y$ splits
  into $x\succ y$ and $x\sim y$, and it needs no completeness: given $x\succeq
  y$, either $y\succeq x$ as well, which is $x\sim y$, or not, which is $x\succ
  y$.
\end{remark}

% ===========================================================================
\exercise{1.B.5}{B}{Representation on a finite set}

\begin{statement}
  Show that if $X$ is finite and $\succeq$ is a rational preference relation on
  $X$, then there is a utility function $u:X\to\mathbb R$ that represents
  $\succeq$. [\emph{Hint}: Consider first the case in which the individual's
  ranking between any two elements of $X$ is strict (i.e., there is never any
  indifference), and construct a utility function representing these preferences;
  then extend your argument to the general case.]
\end{statement}

\begin{nts}
  Exhibit a function $u:X\to\mathbb R$ and show that, for all $x,y\in X$,
  $x\succeq y\iff u(x)\geq u(y)$.
  \medskip

  \noindent
  This is an existence claim, so the proof must produce a candidate. Count the
  alternatives that each $x$ weakly outranks, including those indifferent to it. Write
  \[
  L(x)=\{z\in X: x\succeq z\},
  \qquad
  u(x)=\#L(x),
  \]
  the number of alternatives that $x$ is at least as good as. Since $X$ is
  finite, $u(x)$ is a well-defined integer, so $u:X\to\mathbb R$.
\end{nts}

\begin{proof}
  Use the function $u(x)=\#L(x)$ defined above. Let $x,y\in X$ be arbitrary.
  \begin{pstep}
    $(\Rightarrow)$ Suppose $x\succeq y$.
    \begin{pstep}
      Let $z\in L(y)$ be arbitrary.
      \begin{pstep}
        $y\succeq z$. \by{definition of $L(y)$}

        $x\succeq y$ and $y\succeq z$, so $x\succeq z$. \by{transitivity}

        So $z\in L(x)$. \by{definition of $L(x)$}
      \end{pstep}
      Thus $L(y)\subseteq L(x)$.

      Consequently $u(y)=\#L(y)\leq\#L(x)=u(x)$.
    \end{pstep}
    Thus $x\succeq y\Rightarrow u(x)\geq u(y)$.

    $(\Leftarrow)$ Suppose $u(x)\geq u(y)$.
    \begin{pstep}
      Suppose, for contradiction, that not $x\succeq y$.
      \begin{pstep}
        $y\succeq x$. \by{completeness}

        So $L(x)\subseteq L(y)$.
        \by{forward inclusion, with $x,y$ exchanged}

        $y\succeq y$, so $y\in L(y)$. \by{Step 0 of Exercise 1.B.2}

        $y\notin L(x)$, since $x\succeq y$ fails. \by{definition of $L(x)$}

        Hence $L(x)\subsetneq L(y)$, with $y$ witnessing strictness.

        These sets are finite, so $\#L(x)<\#L(y)$, that is, $u(x)<u(y)$.

        This contradicts the still-active assumption $u(x)\geq u(y)$. \contra
      \end{pstep}
      Thus $x\succeq y$.
    \end{pstep}
    Thus $u(x)\geq u(y)\Rightarrow x\succeq y$.
  \end{pstep}
  Thus, for all $x,y\in X$, $x\succeq y\iff u(x)\geq u(y)$, so $u$ represents
  $\succeq$.
\end{proof}

\begin{remark}
  Finiteness is used in two places: it makes each $\#L(x)$ a real-valued
  integer, and it turns the strict inclusion $L(x)\subsetneq L(y)$ into a
  strict cardinality inequality. For infinite sets a proper subset can have the
  same cardinality as the whole set, and the theorem does not extend to all
  infinite domains: lexicographic preferences on $\mathbb R_{+}^{2}$ are
  complete, transitive, strongly monotone and strictly convex, but admit no
  real-valued utility representation (MWG 3.C.1, and the Day 3 handout).

  One way to follow the hint is to induct on $\#X$ when there is no
  indifference between distinct alternatives, then pass to the indifference
  classes in the general case. That works, but the
  counting function above absorbs indifference with no extra argument, since
  $x\sim y$ gives $L(x)=L(y)$ and hence $u(x)=u(y)$ at once. Both proofs are
  worth having. The induction is the one that generalises to the constructions of
  Day 3; this one is the one to put on a board.
\end{remark}

\newpage
\section*{Choice Rules (Section 1.C)}

In the three-alternative examples, $x,y,z$ denote distinct alternatives.

% ===========================================================================
\exercise{1.C.1}{B}{What the weak axiom leaves open}

\begin{statement}
  Consider the choice structure $(\mathbb B,C(\cdot))$ with
  $\mathbb B=\bigl\{\{x,y\},\{x,y,z\}\bigr\}$ and $C(\{x,y\})=\{x\}$. Show that
  if $(\mathbb B,C(\cdot))$ satisfies the weak axiom, then we must have
  $C(\{x,y,z\})=\{x\}$, $=\{z\}$, or $=\{x,z\}$.
\end{statement}

\begin{nts}
  $C(\{x,y,z\})\in\bigl\{\{x\},\{z\},\{x,z\}\bigr\}$.
  \medskip

  \noindent
  Two facts come free from the definition of a choice rule: $C(\{x,y,z\})$ is
  nonempty, and it is a subset of $\{x,y,z\}$. The nonempty subsets of
  $\{x,y,z\}$ that omit $y$ are exactly $\{x\}$, $\{z\}$ and $\{x,z\}$. So the
  single claim left to prove is
  \[
  y\notin C(\{x,y,z\}).
  \]
\end{nts}

\begin{proof}
  Write $B=\{x,y,z\}$ and $B^{\prime}=\{x,y\}$, both in $\mathbb B$.
  \begin{pstep}
    Suppose, for contradiction, $y\in C(B)$.
    \begin{pstep}
      Apply the weak axiom to the ordered pair $(y,x)$:
      \begin{pstep}
        $y,x\in B$ and $y\in C(B)$. \by{antecedent of Definition 1.C.1}

        $B^{\prime}\in\mathbb B$, $y,x\in B^{\prime}$, and $x\in C(B^{\prime})$.
        \by{$C(B^{\prime})=\{x\}$}

        Therefore $y\in C(B^{\prime})$. \by{weak axiom}
      \end{pstep}
      But $C(B^{\prime})=\{x\}$ and $y\neq x$, so $y\notin C(B^{\prime})$. \contra
    \end{pstep}
    Thus $y\notin C(B)$.
  \end{pstep}
  $C(B)$ is a nonempty subset of $\{x,y,z\}$ that omits $y$, hence a nonempty
  subset of $\{x,z\}$.
  The nonempty subsets of $\{x,z\}$ are $\{x\}$, $\{z\}$ and $\{x,z\}$.
  Thus $C(\{x,y,z\})=\{x\}$, $=\{z\}$, or $=\{x,z\}$.
\end{proof}

\begin{remark}
  The conclusion is tight: each of the three is actually consistent with the weak
  axiom on this domain. The only pair available in both menus is $\{x,y\}$, and
  in each of the three cases $C(\{x,y,z\})\cap\{x,y\}$ is either $\{x\}$ or
  empty, so no reversal is ever exhibited. In particular the weak axiom is silent
  about a comparison with $z$ across these two menus: $z$ is unavailable in
  $\{x,y\}$ and can be the unique choice in $\{x,y,z\}$. Being unavailable is
  not the same as being available but unchosen.

  The step to watch is the application of the axiom. It is stated with a named
  pair $(x,y)$, and here it is used with the names exchanged. Read Definition
  1.C.1 as a schema in two free variables, not as a claim about the letters $x$
  and $y$ in particular.
\end{remark}

% ===========================================================================
\exercise{1.C.2}{B}{A symmetric restatement of the weak axiom}

\begin{statement}
  Show that the weak axiom (Definition 1.C.1) is equivalent to the following
  property holding:
  \begin{quote}
    Suppose that $B,B^{\prime}\in\mathbb B$, that $x,y\in B$, and that
    $x,y\in B^{\prime}$. Then if $x\in C(B)$ and $y\in C(B^{\prime})$, we must have
    $\{x,y\}\subseteq C(B)$ and $\{x,y\}\subseteq C(B^{\prime})$.
  \end{quote}
\end{statement}

\begin{nts}
  Call the displayed property $(P)$. An equivalence is two implications:
  \begin{enumerate}[label=(\alph*),leftmargin=2.4em,topsep=.2em,itemsep=.1em]
    \item WARP $\Rightarrow(P)$;
    \item $(P)\Rightarrow$ WARP.
  \end{enumerate}
  Both go directly. Property $(P)$ displays two conclusions symmetrically,
  whereas Definition 1.C.1 displays only one of them. Thus (a) applies the weak
  axiom twice, exchanging the menus and alternatives; (b) uses just one part of
  the conclusion of $(P)$.
\end{nts}

\begin{proof}[Proof of (a), WARP $\Rightarrow(P)$]
  Suppose the weak axiom holds.
  \begin{pstep}
    Let $B,B^{\prime}\in\mathbb B$ and $x,y\in B\cap B^{\prime}$ be arbitrary.
    \begin{pstep}
      Suppose $x\in C(B)$ and $y\in C(B^{\prime})$.
      \begin{pstep}
        \emph{First application: the pair $(x,y)$ and menus $B,B^{\prime}$.}
        \begin{pstep}
          $x,y\in B$ and $x\in C(B)$.

          $x,y\in B^{\prime}$ and $y\in C(B^{\prime})$.

          Therefore $x\in C(B^{\prime})$. \by{weak axiom}
        \end{pstep}
        \emph{Second application: the pair $(y,x)$ and menus $B^{\prime},B$.}
        \begin{pstep}
          $y,x\in B^{\prime}$ and $y\in C(B^{\prime})$.

          $y,x\in B$ and $x\in C(B)$.

          Therefore $y\in C(B)$. \by{weak axiom}
        \end{pstep}
        Combining these with the assumed memberships gives
        $\{x,y\}\subseteq C(B)$ and $\{x,y\}\subseteq C(B^{\prime})$.
      \end{pstep}
      Thus the hypotheses of $(P)$ imply its conclusion for this choice of menus
      and alternatives.
    \end{pstep}
    Since the menus and alternatives were arbitrary, $(P)$ holds.
  \end{pstep}
  Thus WARP $\Rightarrow(P)$.
\end{proof}

\begin{proof}[Proof of (b), $(P)\Rightarrow$ WARP]
  Suppose $(P)$ holds.
  \begin{pstep}
    Let $B\in\mathbb B$ and $x,y\in B$ be arbitrary.
    \begin{pstep}
      Suppose $x\in C(B)$.
      \begin{pstep}
        Let $B^{\prime}\in\mathbb B$ with $x,y\in B^{\prime}$ and
        $y\in C(B^{\prime})$ be arbitrary.
        \begin{pstep}
          The menus $B,B^{\prime}$ contain both alternatives, and
          $x\in C(B)$ and $y\in C(B^{\prime})$.

          These are the hypotheses of $(P)$, so $\{x,y\}\subseteq C(B^{\prime})$.

          In particular, $x\in C(B^{\prime})$.
        \end{pstep}
        Thus $x\in C(B^{\prime})$ for every such $B^{\prime}$.
      \end{pstep}
      Thus choosing $x$ from $B$ has the implication required by the weak axiom.
    \end{pstep}
    Since $B,x,y$ were arbitrary, the weak axiom holds.
  \end{pstep}
  Thus $(P)\Rightarrow$ WARP.
\end{proof}

\begin{remark}
  $(P)$ is the ``no revealed reversal'' phrasing made precise: if $x$ and $y$ are
  each chosen somewhere in the presence of the other, then wherever both are
  available and either one is chosen, both are. They are chosen as a block or not
  at all --- which is what it means for the data to reveal indifference rather
  than a ranking.
\end{remark}

% ===========================================================================
\exercise{1.C.3}{C}{Two revealed strict preferences}

\begin{statement}
  Suppose that choice structure $(\mathbb B,C(\cdot))$ satisfies the weak axiom.
  Consider the following two possible revealed preferred relations, $\succ^{*}$
  and $\succ^{**}$:
  \[
  \begin{aligned}
    x\succ^{*}y \quad&\Longleftrightarrow\quad
    \text{there is some }B\in\mathbb B\text{ such that }x,y\in B,\ x\in C(B),
    \text{ and } y\notin C(B)\\
    x\succ^{**}y \quad&\Longleftrightarrow\quad
    x\succeq^{*}y\text{ but not }y\succeq^{*}x
  \end{aligned}
  \]
  where $\succeq^{*}$ is the revealed at-least-as-good-as relation defined in
  Definition 1.C.2.
  \begin{enumerate}[label=(\alph*),leftmargin=2.4em,topsep=.3em,itemsep=.25em]
    \item Show that $\succ^{*}$ and $\succ^{**}$ give the same relation over $X$;
    that is, for any $x,y\in X$, $x\succ^{*}y\iff x\succ^{**}y$. Is this
    still true if $(\mathbb B,C(\cdot))$ does not satisfy the weak axiom?
    \item Must $\succ^{*}$ be transitive?
    \item Show that if $\mathbb B$ includes all three-element subsets of $X$,
    then $\succ^{*}$ is transitive.
  \end{enumerate}
\end{statement}

\begin{nts}
  \begin{enumerate}[label=(\alph*),leftmargin=2.4em,topsep=.2em,itemsep=.35em]
    \item For all $x,y\in X$, two implications:
    $x\succ^{*}y\Rightarrow x\succ^{**}y$ and
    $x\succ^{**}y\Rightarrow x\succ^{*}y$. Then, for the second question,
    either a proof that the equivalence survives without the weak axiom, or
    a choice structure violating the axiom in which the two relations
    differ.
    \item To answer no, exhibit a choice structure satisfying the weak axiom
    and alternatives $x,y,z$ with $x\succ^{*}y$, $y\succ^{*}z$ and not
    $x\succ^{*}z$. Both the axiom and the failure of transitivity must be
    checked.
    \item For all $x,y,z\in X$: if $x\succ^{*}y$ and $y\succ^{*}z$, then
    $x\succ^{*}z$.
  \end{enumerate}
  Unfold once, so that the definitions are in view. $x\succ^{*}y$ asserts the
  existence of a menu containing both $x$ and $y$ on which $x$ is chosen and
  $y$ is not. $x\succeq^{*}y$
  asserts the existence of a menu containing both on which $x$ is chosen. So
  $\succ^{**}$ is the asymmetric part of $\succeq^{*}$: $x$ is revealed at least
  as good as $y$, and $y$ is never revealed at least as good as $x$.
\end{nts}

\begin{proof}[Proof of (a)]
  Let $x,y\in X$ be arbitrary.
  \begin{pstep}
    $(\Rightarrow)$ Suppose $x\succ^{*}y$.
    \begin{pstep}
      Choose a witnessing menu $B\in\mathbb B$: $x,y\in B$, $x\in C(B)$,
      and $y\notin C(B)$.

      The same $B$ witnesses $x\succeq^{*}y$. \by{Definition 1.C.2}

      Suppose, for contradiction, $y\succeq^{*}x$.
      \begin{pstep}
        Choose $B^{\prime}\in\mathbb B$ with $x,y\in B^{\prime}$ and
        $y\in C(B^{\prime})$. \by{Definition 1.C.2}

        Apply WARP to $(y,x)$, first on $B^{\prime}$ and then on $B$:
        $y$ is chosen on $B^{\prime}$, and $x$ is chosen on $B$, with both
        alternatives available in both menus.

        Therefore $y\in C(B)$. \by{weak axiom}

        But $y\notin C(B)$ by the choice of $B$. \contra
      \end{pstep}
      Thus not $y\succeq^{*}x$.

      Together with $x\succeq^{*}y$, this gives $x\succ^{**}y$.
    \end{pstep}
    Thus $x\succ^{*}y\Rightarrow x\succ^{**}y$.

    $(\Leftarrow)$ Suppose $x\succ^{**}y$.
    \begin{pstep}
      Then $x\succeq^{*}y$ and not $y\succeq^{*}x$.
      \by{definition of $\succ^{**}$}

      Choose $B\in\mathbb B$ with $x,y\in B$ and $x\in C(B)$.
      \by{$x\succeq^{*}y$}

      Suppose, for contradiction, $y\in C(B)$.
      \begin{pstep}
        Then $B$ witnesses $y\succeq^{*}x$, since $y,x\in B$ and $y\in C(B)$.

        This contradicts not $y\succeq^{*}x$. \contra
      \end{pstep}
      Thus $y\notin C(B)$.

      So $B$ witnesses $x\succ^{*}y$: $x,y\in B$, $x\in C(B)$, $y\notin C(B)$.
    \end{pstep}
    Thus $x\succ^{**}y\Rightarrow x\succ^{*}y$.
  \end{pstep}
  Thus, for all $x,y\in X$, $x\succ^{*}y\iff x\succ^{**}y$.
\end{proof}

\noindent
\emph{Without the weak axiom the equivalence fails.} The direction
$(\Leftarrow)$ used no weak axiom and so survives; $(\Rightarrow)$ does not.
Take
\[
X=\{x,y,z\},\qquad
\mathbb B=\bigl\{\{x,y\},\{x,y,z\}\bigr\},\qquad
C(\{x,y\})=\{x\},\quad C(\{x,y,z\})=\{y\}.
\]
Then $\{x,y\}$ witnesses $x\succ^{*}y$, and $\{x,y,z\}$ witnesses
$y\succ^{*}x$. But $x\succeq^{*}y$ and $y\succeq^{*}x$ both hold, so neither
$x\succ^{**}y$ nor $y\succ^{**}x$. The two relations differ. This structure
violates the weak axiom, as Exercise 1.C.1 shows it must.

\begin{proof}[Proof of (b): no]
  Take
  \[
  \begin{gathered}
    X=\{x,y,z\},\qquad
    \mathbb B=\bigl\{\{x,y\},\{y,z\}\bigr\},\\
    C(\{x,y\})=\{x\},\qquad C(\{y,z\})=\{y\}.
  \end{gathered}
  \]
  \emph{First check the weak axiom.}
  \begin{pstep}
    Within a single menu, the required conclusion is already a premise.

    The two different menus share only $y$, so no pair of distinct alternatives
    is available in both. For a pair with equal alternatives, the required
    conclusion is again already a premise.
  \end{pstep}
  Thus this choice structure satisfies the weak axiom.

  \emph{Now check the failure of transitivity.}
  \begin{pstep}
    $x\succ^{*}y$, witnessed by $\{x,y\}$.

    $y\succ^{*}z$, witnessed by $\{y,z\}$.

    No menu in $\mathbb B$ contains both $x$ and $z$, so no menu can witness
    $x\succ^{*}z$.

    Thus not $x\succ^{*}z$.
  \end{pstep}
  Therefore $\succ^{*}$ need not be transitive, even under the weak axiom.
\end{proof}

\begin{remark}
  This is the point the rich-domain slide is making. The failure is not a failure
  of the axiom, which holds here; it is a failure of the data. The comparison of
  $x$ against $z$ was never run, so nothing could have revealed it. Transitivity
  is a property one can only observe on a domain wide enough to contain the test.
\end{remark}

\begin{proof}[Proof of (c)]
  Assume $\mathbb B$ contains every three-element subset of $X$. The weak axiom
  continues to hold by the standing assumption of this exercise.
  \begin{pstep}
    Let $x,y,z\in X$ be arbitrary.
    \begin{pstep}
      Suppose $x\succ^{*}y$ and $y\succ^{*}z$.
      \begin{pstep}
        By (a), $\succ^{*}=\succ^{**}$. For any $a,b$, $a\succ^{**}b$ requires
        not $b\succeq^{*}a$, whereas $b\succ^{**}a$ would require $b\succeq^{*}a$.
        Thus $\succ^{*}$ is asymmetric and hence irreflexive.

        Consequently $x\neq y$ and $y\neq z$.

        Suppose, for contradiction, $x=z$.
        \begin{pstep}
          Then the two hypotheses give $x\succ^{*}y$ and $y\succ^{*}x$,
          contradicting asymmetry. \contra
        \end{pstep}
        Thus $x\neq z$, so $x,y,z$ are three distinct alternatives.

        Put $T=\{x,y,z\}$. Then $T\in\mathbb B$.
        \by{the domain contains every three-element subset}

        By (a) and the two strict comparisons, not $y\succeq^{*}x$ and
        not $z\succeq^{*}y$.

        Suppose, for contradiction, $y\in C(T)$.
        \begin{pstep}
          Then $T$ witnesses $y\succeq^{*}x$, since $y,x\in T$.

          This contradicts not $y\succeq^{*}x$. \contra
        \end{pstep}
        Thus $y\notin C(T)$.

        Suppose, for contradiction, $z\in C(T)$.
        \begin{pstep}
          Then $T$ witnesses $z\succeq^{*}y$, since $z,y\in T$.

          This contradicts not $z\succeq^{*}y$. \contra
        \end{pstep}
        Thus $z\notin C(T)$.

        Since $C(T)\subseteq T$, we have $C(T)\subseteq\{x\}$.

        Also $C(T)\neq\varnothing$, so $C(T)=\{x\}$.
        \by{nonemptiness of choice}

        Hence $T$ witnesses $x\succ^{*}z$: $x,z\in T$, $x\in C(T)$, and
        $z\notin C(T)$.
      \end{pstep}
      Thus $x\succ^{*}y$ and $y\succ^{*}z$ imply $x\succ^{*}z$.
    \end{pstep}
    Since $x,y,z$ were arbitrary, $\succ^{*}$ is transitive.
  \end{pstep}
  Thus, under WARP, the inclusion of every three-element menu implies
  transitivity of $\succ^{*}$.
\end{proof}

\begin{remark}
  Compare with part (b): the triple menu makes the comparison of $x$ with $z$
  observable, and nonemptiness forces $C(T)=\{x\}$. This is the triple-menu
  argument behind transitivity in the Day 1 theorem. Here we proved
  transitivity of $\succ^{*}$, not completeness of revealed preference.
\end{remark}

\end{document}
